\begin{definition}[Open middle tensor at an edge]%
    \label{def:peps_edgeOpenMiddleWeight}
    \lean{TNLean.PEPS.edgeOpenMiddleWeight}
    \leanok
    \uses{def:peps_edgeInsertedBoundaryConfig, def:peps_edgeMiddleVertices}
    Fixing the residual endpoint data but leaving the bond $e$ out of the
    middle region, the open middle tensor is the contraction over all virtual
    indices away from $e$ that are compatible with those endpoint residual
    data.
\end{definition}

\begin{definition}[Middle physical configurations at an edge]%
    \label{def:peps_edgeMiddlePhysicalConfig}
    \lean{TNLean.PEPS.EdgeMiddlePhysicalConfig}
    \lean{TNLean.PEPS.edgeMiddlePhysicalConfigOf}
    \leanok
    \uses{def:peps_edgeMiddleVertices}
    For an edge $e=(u,v)$, the middle physical configurations are the
    physical indices on the vertices in $V\setminus\{u,v\}$. A global
    physical configuration restricts to such a middle configuration.
\end{definition}

\begin{definition}[Middle tensor with middle physical index]%
    \label{def:peps_edgeMiddleWeightOn}
    \lean{TNLean.PEPS.edgeMiddleWeightOn}
    \lean{TNLean.PEPS.edgeOpenMiddleWeightOn}
    \leanok
    \uses{def:peps_edgeMiddlePhysicalConfig,
        def:peps_edgeMiddleConfig, def:peps_edgeOpenMiddleWeight}
    The ordinary and open middle contractions may be written using only the
    middle physical configuration, rather than a physical configuration on all
    vertices. This is the middle tensor in the three-site chain obtained after
    the edge blocking in arXiv:1804.04964, Section~3.
\end{definition}

\begin{theorem}[Middle restriction of the blocked tensor]%
    \label{thm:peps_edgeMiddleWeight_eq_on}
    \lean{TNLean.PEPS.edgeMiddleWeight_eq_on}
    \lean{TNLean.PEPS.edgeOpenMiddleWeight_eq_on}
    \leanok
    \uses{def:peps_edgeMiddleWeightOn,
        def:peps_edgeMiddlePhysicalConfig}
    The middle contractions written with a global physical configuration are
    the corresponding middle-indexed contractions evaluated on the restricted
    middle physical configuration.
\end{theorem}
\begin{proof}\leanok
    Rewrite the product over $V\setminus\{u,v\}$ as the product over the
    corresponding finite set of middle vertices.
\end{proof}

\begin{definition}[Restricting middle configurations away from the edge]%
    \label{def:peps_edgeMiddleConfigToOpenMiddleConfig}
    \lean{TNLean.PEPS.edgeMiddleConfigToOpenMiddleConfig}
    \leanok
    \uses{def:peps_edgeMiddleConfig, def:peps_edgeOpenMiddleWeight}
    An ordinary middle configuration for the edge-blocked contraction restricts
    to a virtual configuration on all edges except the distinguished edge. This
    is the first reindexing map needed to compare the identity insertion with
    the ordinary edge-blocked middle tensor.
\end{definition}

\begin{definition}[Restoring the distinguished edge index]%
    \label{def:peps_edgeOpenMiddleConfigToMiddleConfig}
    \lean{TNLean.PEPS.edgeOpenMiddleConfigToMiddleConfig}
    \leanok
    \uses{def:peps_edgeBoundaryConfig, def:peps_edgeMiddleConfigToOpenMiddleConfig}
    Conversely, a compatible open middle configuration becomes an ordinary
    middle configuration after the fixed edge-boundary index is restored on the
    distinguished edge. Together with the restriction map, this gives the
    finite reindexing underlying the identity-insertion specialization.
\end{definition}

\begin{definition}[Open-middle reindexing equivalence]%
    \label{def:peps_edgeMiddleConfigEquivOpenMiddleConfig}
    \lean{TNLean.PEPS.edgeMiddleConfigEquivOpenMiddleConfig}
    \leanok
    \uses{def:peps_edgeMiddleConfigToOpenMiddleConfig,
        def:peps_edgeOpenMiddleConfigToMiddleConfig}
    For fixed edge-boundary data, ordinary middle configurations are equivalent
    to compatible open-middle configurations. This is the finite reindexing
    used when the inserted matrix is the identity.
\end{definition}

\begin{theorem}[Middle-indexed open tensor reindexing]%
    \label{thm:peps_edgeMiddleWeightOn_eq_edgeOpenMiddleWeightOn}
    \lean{TNLean.PEPS.edgeMiddleWeightOn_eq_edgeOpenMiddleWeightOn}
    \leanok
    \uses{def:peps_edgeMiddleWeightOn,
        def:peps_edgeMiddleConfigEquivOpenMiddleConfig}
    The ordinary middle tensor and the open middle tensor agree after restoring
    the fixed distinguished-edge index and reindexing the finite sum, with the
    physical index already restricted to the middle block.
\end{theorem}
\begin{proof}\leanok
    Reindex the ordinary middle configurations by the equivalence with
    compatible open-middle configurations, then compare each local tensor entry.
\end{proof}

\begin{definition}[Edge insertion coefficient]%
    \label{def:peps_edgeInsertedCoeff}
    \lean{TNLean.PEPS.edgeInsertedCoeff}
    \leanok
    \uses{def:peps_edgeInsertedBoundaryConfig, def:peps_edgeOpenMiddleWeight}
    For an edge $e=(u,v)$ and a matrix $X$ acting on the bond space of $e$,
    the edge insertion coefficient is the contraction obtained by inserting
    $X$ between the endpoint tensors and the open middle tensor:
    \begin{center}
        \TNPEPSEdgeInsertedCoeff
    \end{center}
\end{definition}

\begin{theorem}[Middle tensor reindexing for an opened edge]%
    \label{thm:peps_edgeMiddleWeight_eq_edgeOpenMiddleWeight}
    \lean{TNLean.PEPS.edgeMiddleWeight_eq_edgeOpenMiddleWeight}
    \leanok
    \uses{def:peps_edgeOpenMiddleWeight,
        def:peps_edgeMiddleConfigEquivOpenMiddleConfig}
    Restoring the fixed distinguished-edge index identifies the ordinary blocked
    middle tensor with the open middle tensor. This is the finite reindexing
    used before the identity matrix collapses the two inserted bond indices to
    the diagonal.
\end{theorem}

\begin{theorem}[Diagonal summand of the identity insertion]%
    \label{thm:peps_edgeInsertedCoeff_identity_diagonal_summand}
    \lean{TNLean.PEPS.edgeInsertedCoeff_identity_diagonal_summand}
    \leanok
    \uses{def:peps_edgeBoundaryToInsertedBoundaryConfig,
        thm:peps_edgeMiddleWeight_eq_edgeOpenMiddleWeight}
    On the diagonal boundary datum selected by the identity matrix, the
    inserted-edge summand is the ordinary edge-blocked summand. Thus the
    remaining identity-insertion step is purely a finite summation statement:
    the off-diagonal inserted boundary data vanish and the diagonal data reindex
    the ordinary edge-boundary sum.
\end{theorem}

\begin{theorem}[Off-diagonal summands of the identity insertion]%
    \label{thm:peps_edgeInsertedCoeff_identity_offDiagonal_summand}
    \lean{TNLean.PEPS.edgeInsertedCoeff_identity_offDiagonal_summand}
    \leanok
    \uses{def:peps_edgeInsertedCoeff}
    If the two inserted bond indices are distinct, then the corresponding
    summand in the identity insertion is zero. This is the pointwise
    off-diagonal vanishing used before the inserted-boundary sum is restricted
    to its diagonal part.
\end{theorem}

\begin{theorem}[Identity insertion recovers the edge-blocked coefficient]%
    \label{thm:peps_edgeInsertedCoeff_identity}
    \lean{TNLean.PEPS.edgeInsertedCoeff_identity}
    \leanok
    \uses{def:peps_edgeInsertedCoeff, def:peps_edgeBlockedCoeff,
        def:peps_edgeBoundaryToInsertedBoundaryConfig,
        def:peps_edgeBoundaryConfigEquivDiagonalInsertedBoundaryConfig,
        thm:peps_edgeInsertedCoeff_identity_diagonal_summand,
        thm:peps_edgeInsertedCoeff_identity_offDiagonal_summand}
    If the inserted matrix on the distinguished edge is the identity, then the
    edge insertion coefficient is the ordinary edge-blocked coefficient.
\end{theorem}
\begin{proof}\leanok
    Split the inserted-boundary sum into the diagonal part, where the two
    inserted bond indices agree, and the off-diagonal part. The off-diagonal
    summands vanish for the identity matrix, while the diagonal summands
    reindex through the equivalence with ordinary edge-boundary data.
\end{proof}

\begin{definition}[Blocked middle tensor at an edge]%
    \label{def:peps_edgeMiddleWeight}
    \lean{TNLean.PEPS.edgeMiddleWeight}
    \leanok
    \uses{def:peps_edgeBoundaryMatches, def:peps_edgeMiddleConfig,
        def:peps_edgeMiddleVertices}
    Fixing the edge-centered boundary data, the blocked middle tensor is the sum
    over global virtual configurations compatible with those data of the product
    of local tensor entries over the middle region $V\setminus\{u,v\}$.
    The graph-local contraction is read as a three-site chain:
    \begin{center}
        \TNPEPSEdgeBlockingReduction
    \end{center}
\end{definition}

\begin{definition}[Edge-blocked PEPS coefficient]%
    \label{def:peps_edgeBlockedCoeff}
    \lean{TNLean.PEPS.edgeBlockedCoeff}
    \leanok
    \uses{def:peps_edgeLeftLocalConfig, def:peps_edgeMiddleWeight,
        def:peps_edgeRightLocalConfig}
    The edge-blocked coefficient is the three-block contraction obtained from
    the left endpoint tensor, the blocked middle tensor, and the right endpoint
    tensor at the chosen edge.
\end{definition}

\begin{theorem}[Edge-blocked coefficient equals the PEPS coefficient]%
    \label{thm:peps_edgeBlockedCoeff_eq_stateCoeff}
    \lean{TNLean.PEPS.edgeBlockedCoeff_eq_stateCoeff}
    \leanok
    \uses{def:peps_edgeBlockedCoeff, thm:peps_stateCoeff_splitAtEdge,
        def:peps_virtualConfigEquivEdgeBoundary}
    For every edge $e$, the edge-blocked three-block coefficient is exactly the
    original PEPS state coefficient.
\end{theorem}
\begin{proof}\leanok
    \uses{def:peps_virtualConfigEquivEdgeBoundary,
        thm:peps_edgeLeftLocalConfig_eq_of_boundaryMatches,
        thm:peps_edgeRightLocalConfig_eq_of_boundaryMatches}
    Write $\operatorname{Coeff}_A(\sigma)$ for the PEPS state coefficient and
    $C_A(e;\sigma)$ for the edge-blocked coefficient, and put $e=(u,v)$ and
    $M_e=V\setminus\{u,v\}$. The equivalence
    $\omega\leftrightarrow(\beta,\omega_{M_e})$ between global virtual
    configurations and an edge boundary configuration with a compatible middle
    configuration rewrites the PEPS coefficient as
    \[
        \sum_{\omega}
        A_u(\omega|_u,\sigma_u)
        \left(\prod_{x\in M_e}A_x(\omega|_x,\sigma_x)\right)
        A_v(\omega|_v,\sigma_v)
        =
        \sum_{\beta}
        A_u(\beta_u,\sigma_u)
        C_{A,M_e}(\beta,\sigma|_{M_e})
        A_v(\beta_v,\sigma_v).
    \]
    Under the equivalence, $\omega|_u=\beta_u$ and $\omega|_v=\beta_v$; the
    inner sum over $\omega_{M_e}$ is exactly $C_{A,M_e}(\beta,\sigma|_{M_e})$.
    Thus $C_A(e;\sigma)=\operatorname{Coeff}_A(\sigma)$.
\end{proof}

\begin{theorem}[Identity insertion equals the PEPS coefficient]%
    \label{thm:peps_edgeInsertedCoeff_identity_eq_stateCoeff}
    \lean{TNLean.PEPS.edgeInsertedCoeff_identity_eq_stateCoeff}
    \leanok
    \uses{thm:peps_edgeInsertedCoeff_identity,
        thm:peps_edgeBlockedCoeff_eq_stateCoeff}
    If the inserted matrix on the distinguished edge is the identity, then the
    edge insertion coefficient is the original PEPS state coefficient.
\end{theorem}
\begin{proof}\leanok
    Let $e=(u,v)$ and put $M_e=V\setminus\{u,v\}$. With $\Id$ inserted on $e$, the
    inserted-boundary sum is
    \[
        C_A(e,\Id;\sigma)
        =
        \sum_{a,b,\lambda,\rho}
        \delta_{a,b}
        A_u(a,\lambda;\sigma_u)
        C^{\mathrm{open}}_{A,M_e}(\lambda,\rho;\sigma|_{M_e})
        A_v(b,\rho;\sigma_v).
    \]
    Here $a,b$ are the two copies of the distinguished bond index, while
    $\lambda,\rho$ are the remaining virtual configurations at $u$ and $v$.
    The factor $\delta_{a,b}$
    leaves only the diagonal $a=b$, and the open-middle reindexing gives
    $C^{\mathrm{open}}_{A,M_e}(\lambda,\rho;\sigma|_{M_e})
        =C_{A,M_e}(a,\lambda,\rho;\sigma|_{M_e})$. Hence
    $C_A(e,\Id;\sigma)=C_A(e;\sigma)=\operatorname{Coeff}_A(\sigma)$.
\end{proof}

\begin{theorem}[Same state gives the same edge-blocked coefficients]%
    \label{thm:peps_sameState_edgeBlockedCoeff_eq}
    \lean{TNLean.PEPS.SameState.edgeBlockedCoeff_eq}
    \leanok
    \uses{def:peps_same_state, thm:peps_edgeBlockedCoeff_eq_stateCoeff}
    If two PEPS tensors have the same state, then their edge-blocked
    coefficients agree at every edge and every physical configuration.
\end{theorem}
\begin{proof}\leanok
    For every physical configuration $\sigma$,
    \[
        C_A(e;\sigma)
        =
        \operatorname{Coeff}_A(\sigma)
        =
        \operatorname{Coeff}_B(\sigma)
        =
        C_B(e;\sigma).
    \]
\end{proof}

\begin{theorem}[Same state gives the same identity insertions]%
    \label{thm:peps_sameState_edgeInsertedCoeff_identity_eq}
    \lean{TNLean.PEPS.SameState.edgeInsertedCoeff_identity_eq}
    \leanok
    \uses{def:peps_same_state,
        thm:peps_edgeInsertedCoeff_identity_eq_stateCoeff}
    If two PEPS tensors have the same state, then their edge-centered identity
    insertions agree at every edge and every physical configuration.
\end{theorem}
\begin{proof}\leanok
    For every physical configuration $\sigma$,
    \[
        C_A(e,\Id;\sigma)
        =
        \operatorname{Coeff}_A(\sigma)
        =
        \operatorname{Coeff}_B(\sigma)
        =
        C_B(e,\Id;\sigma).
    \]
\end{proof}

\begin{definition}[Middle tensor family at an edge]%
    \label{def:peps_edgeMiddleTensorFamily}
    \lean{TNLean.PEPS.edgeMiddleTensorFamily}
    \lean{TNLean.PEPS.EdgeMiddleTensorInjective}
    \leanok
    \uses{def:peps_edgeMiddleWeightOn}
    For the edge-blocked three-site chain at $e=(u,v)$, the middle tensor is
    the family obtained by fixing the two residual endpoint boundaries and
    evaluating the open middle contraction on the middle physical index.
    Its injectivity is the middle-block part of the assertion following
    the edge-blocking diagram.
\end{definition}

\begin{lemma}[Middle tensor vanishes for an empty complement configuration]%
    \label{lem:peps_edgeMiddleTensorFamily_eq_zero_of_isEmpty}
    \lean{TNLean.PEPS.edgeMiddleTensorFamily_eq_zero_of_isEmpty}
    \leanok
    \uses{def:peps_edgeMiddleTensorFamily}
    If the virtual configuration space on the complement of the distinguished
    edge is empty, then for every residual boundary label $\rho$ the vector
    $T_{A,M_e}^{\rho}$ in the edge-middle tensor family is zero.
\end{lemma}
\begin{proof}\leanok
    The finite sum defining the open middle tensor is taken over the empty
    complement configuration space, and is therefore zero.
\end{proof}

\begin{lemma}[Failure of middle injectivity for an empty complement configuration]%
    \label{lem:peps_not_edgeMiddleTensorInjective_of_isEmpty}
    \lean{TNLean.PEPS.not_edgeMiddleTensorInjective_of_isEmpty}
    \leanok
    \uses{def:peps_edgeMiddleTensorFamily,
        lem:peps_edgeMiddleTensorFamily_eq_zero_of_isEmpty}
    If the complement configuration space is empty and the residual
    boundary-label set is nonempty, then the edge-middle tensor family is not
    injective.
\end{lemma}
\begin{proof}\leanok
    Choose a residual boundary label. By the preceding lemma its middle-tensor
    vector is the zero vector. A linearly independent family cannot contain a
    zero vector.
\end{proof}

\begin{definition}[Edge-blocked three-site injectivity]%
    \label{def:peps_edgeBlockedThreeSiteInjective}
    \lean{TNLean.PEPS.EdgeBlockedThreeSiteInjective}
    \leanok
    \uses{def:peps_localTensorMap, def:peps_edgeMiddleTensorFamily}
    For a chosen edge $e=(u,v)$, the edge-blocked three-site chain is
    injective when the left endpoint tensor map, the middle tensor family, and
    the right endpoint tensor map are all injective.
\end{definition}

\begin{definition}[Singleton blocked region]%
    \label{def:peps_singletonRegionTensor}
    \lean{TNLean.PEPS.SingletonRegionPhysicalConfig,
        TNLean.PEPS.singletonRegionTensorFamily,
        TNLean.PEPS.SingletonRegionTensorInjective,
        TNLean.PEPS.SingletonRegionInjectivityComparison}
    \leanok
    \uses{def:peps_regionInjectivityData}
    For a singleton block $\{v\}$, the blocked physical index is the original
    physical index and the blocked tensor is
    \[
        T_{A,\{v\}}(\eta,\sigma)=A_v(\eta,\sigma).
    \]
    Its injectivity is $\operatorname{inj}(T_{A,\{v\}})$. The comparison with a
    region-injectivity predicate $\kappa$ records the implication
    \[
        \operatorname{inj}(T_{A,\{v\}})\Longrightarrow\kappa(\{v\}).
    \]
\end{definition}

\begin{theorem}[Singleton comparison for a vertex-injective tensor]%
    \label{thm:peps_singletonRegionInjectivityComparison_of_isVertexInjective}
    \lean{TNLean.PEPS.singletonRegionInjectivityComparison_of_isVertexInjective}
    \leanok
    \uses{def:IsVertexInjective, def:peps_singletonRegionTensor,
        def:peps_regionInjectivityDataOf,
        thm:peps_regionBlockedTensorInjective_of_isVertexInjective}
    Let $A$ be vertex-injective with every bond dimension positive. Then the
    concrete region-injectivity predicate of $A$ satisfies the singleton
    comparison: injectivity of a singleton blocked tensor implies injectivity of
    the corresponding singleton region.
\end{theorem}
\begin{proof}\leanok
    Under vertex injectivity and positive bond dimensions every finite region
    is injective. In particular the singleton region $\{v\}$ is injective, so
    the comparison holds independently of the singleton-tensor hypothesis.
\end{proof}

\begin{theorem}[Vertex injectivity gives singleton-block injectivity]%
    \label{thm:peps_singletonRegionTensorInjective}
    \lean{TNLean.PEPS.IsVertexInjective.singletonRegionTensorInjective}
    \leanok
    \uses{def:IsVertexInjective, def:peps_singletonRegionTensor}
    If $A$ is vertex-injective, then every singleton blocked tensor is
    injective. For every vertex $v$,
    \[
        \operatorname{inj}(T_{A,\{v\}}).
    \]
\end{theorem}
\begin{proof}\leanok
    For $R=\{v\}$, the contraction defining $T_{A,R}$ has one factor. Hence
    \[
        T_{A,\{v\}}(\eta,\sigma)=A_v(\eta,\sigma).
    \]
    Thus $\operatorname{inj}(A_v)$ is exactly
    $\operatorname{inj}(T_{A,\{v\}})$. Vertex injectivity gives
    $\forall v,\operatorname{inj}(A_v)$, and therefore
    $\forall v,\operatorname{inj}(T_{A,\{v\}})$. For a fixed $v$, this is the
    displayed assertion.
\end{proof}

\begin{theorem}[Singleton comparison gives the region base case]%
    \label{thm:peps_singletonRegionInjectivityFromVertexInjectivity}
    \lean{TNLean.PEPS.SingletonRegionInjectivityFromVertexInjectivity,
        TNLean.PEPS.SingletonRegionInjectivityFromVertexInjectivity.ofSingletonComparison}
    \leanok
    \uses{def:IsVertexInjective, def:peps_singletonRegionTensor,
        thm:peps_singletonRegionTensorInjective}
    Assume that a region-injectivity predicate $\kappa$ satisfies, for every
    vertex $v$,
    \[
        \operatorname{inj}(T_{A,\{v\}})\Longrightarrow\kappa(\{v\}).
    \]
    Then, for every vertex $v$,
    \[
        A\text{ vertex-injective}\Longrightarrow\kappa(\{v\}).
    \]
\end{theorem}
\begin{proof}\leanok
    For each $v$, the preceding theorem gives
    $A\text{ vertex-injective}\Longrightarrow
        \operatorname{inj}(T_{A,\{v\}})$. The comparison hypothesis gives
    $\operatorname{inj}(T_{A,\{v\}})\Longrightarrow\kappa(\{v\})$. Their
    composition is the claimed implication.
\end{proof}

\begin{theorem}[Finite regions from singleton regions and union closure]%
    \label{thm:peps_finiteRegionInjectivityFromSingletons}
    \lean{TNLean.PEPS.RegionInjectivityUnionClosure.finset_injective_of_singletons}
    \leanok
    \uses{thm:peps_singletonRegionInjectivityFromVertexInjectivity,
        lem:peps_regionInjectivityUnionClosure_finite}
    Assume that $\kappa$ satisfies singleton comparison and union closure:
    $\operatorname{inj}(T_{A,\{v\}})\Longrightarrow\kappa(\{v\})$ for every
    vertex $v$, and
    $\kappa(R)\wedge\kappa(S)\Longrightarrow\kappa(R\cup S)$. If $A$ is
    vertex-injective, then every nonempty finite region $R$ satisfies
    $\kappa(R)$.
\end{theorem}
\begin{proof}\leanok
    For $R\ne\varnothing$, write $R=\bigcup_{v\in R}\{v\}$. The preceding
    theorem turns singleton comparison and vertex injectivity into
    $\kappa(\{v\})$ for every $v\in R$. Applying the finite-union consequence
    of union closure to these singleton regions gives $\kappa(R)$.
\end{proof}

\begin{definition}[Region injectivity supplies the edge-middle tensor]%
    \label{def:peps_edgeMiddleRegionInjectivityComparison}
    \lean{TNLean.PEPS.EdgeMiddleRegionInjectivityComparison}
    \leanok
    \uses{def:peps_regionInjectivityData,
        def:peps_edgeMiddleTensorFamily}
    A comparison from finite-region injectivity to the edge-middle tensor
    records the following assertion: if the middle vertex region
    $V\setminus\{u,v\}$ is injective after blocking, then the middle tensor
    in the edge-blocked three-site chain is injective. This isolates the
    finite-region contraction claim in
    \cite[Section~3]{Molnar2018NormalPEPS}.
\end{definition}

\begin{theorem}[Endpoint injectivity in the edge-blocked chain]%
    \label{thm:peps_edgeBlockedEndpointTensorMaps_injective}
    \lean{TNLean.PEPS.IsVertexInjective.edgeBlockedEndpointTensorMaps_injective}
    \leanok
    \uses{def:IsVertexInjective}
    If the original PEPS tensor is vertex-injective, then the two endpoint
    tensor maps in the edge-blocked three-site chain are injective.
\end{theorem}
\begin{proof}\leanok
    If $e=(u,v)$, vertex injectivity gives
    $\operatorname{inj}(T_{A,u})$ and $\operatorname{inj}(T_{A,v})$, which are
    the left and right endpoint tensor maps of the edge-blocked chain.
\end{proof}

\begin{theorem}[Endpoint injectivity at all edges]%
    \label{thm:peps_edgeBlockedEndpointTensorMaps_injective_all}
    \lean{TNLean.PEPS.IsVertexInjective.edgeBlockedEndpointTensorMaps_injective_all}
    \leanok
    \uses{def:IsVertexInjective,
        thm:peps_edgeBlockedEndpointTensorMaps_injective}
    If the original PEPS tensor is vertex-injective, then the two endpoint
    tensor maps in the edge-blocked three-site chain are injective at every
    edge.
\end{theorem}
\begin{proof}\leanok
    For each $e=(u,v)$, Theorem~\ref{thm:peps_edgeBlockedEndpointTensorMaps_injective}
    gives
    $\operatorname{inj}(T_{A,u})\wedge\operatorname{inj}(T_{A,v})$.
\end{proof}

\begin{theorem}[Reducing edge-blocked injectivity to the middle block]%
    \label{thm:peps_edgeBlockedThreeSiteInjective_of_middle}
    \lean{TNLean.PEPS.IsVertexInjective.edgeBlockedThreeSiteInjective_of_middle}
    \leanok
    \uses{def:IsVertexInjective, def:peps_edgeBlockedThreeSiteInjective,
        def:peps_edgeMiddleTensorFamily}
    For a vertex-injective PEPS tensor, injectivity of the edge-blocked middle
    tensor gives the full injectivity statement for the edge-blocked three-site
    chain.
\end{theorem}
\begin{proof}\leanok
    Let $e=(u,v)$ and put $M_e=V\setminus\{u,v\}$. The endpoint theorem supplies
    $\operatorname{inj}(T_{A,u})$ and $\operatorname{inj}(T_{A,v})$, while the
    hypothesis is $\operatorname{inj}(T_{A,M_e})$. Hence
    \[
        \operatorname{inj}(T_{A,u})
        \wedge \operatorname{inj}(T_{A,M_e})
        \wedge \operatorname{inj}(T_{A,v}).
    \]
\end{proof}

\begin{theorem}[All edge-blocked chains from middle-tensor injectivity]%
    \label{thm:peps_edgeBlockedThreeSiteInjective_all_of_middle}
    \lean{TNLean.PEPS.IsVertexInjective.edgeBlockedThreeSiteInjective_all_of_middle}
    \leanok
    \uses{def:IsVertexInjective, def:peps_edgeBlockedThreeSiteInjective,
        def:peps_edgeMiddleTensorFamily,
        thm:peps_edgeBlockedThreeSiteInjective_of_middle}
    If every edge-middle tensor is injective, then a vertex-injective PEPS
    tensor gives an injective edge-blocked three-site chain at every edge.
\end{theorem}
\begin{proof}\leanok
    For each $e=(u,v)$, put $M_e=V\setminus\{u,v\}$. Vertex injectivity gives
    $\operatorname{inj}(T_{A,u})$ and $\operatorname{inj}(T_{A,v})$, while the
    hypothesis gives $\operatorname{inj}(T_{A,M_e})$. Hence
    \[
        \operatorname{inj}(T_{A,u})\wedge
        \operatorname{inj}(T_{A,M_e})\wedge
        \operatorname{inj}(T_{A,v})
    \]
    is the injectivity assertion for the edge-blocked three-site chain at $e$.
\end{proof}

\begin{theorem}[Region injectivity gives edge-blocked injectivity]%
    \label{thm:peps_edgeBlockedThreeSiteInjective_of_region}
    \lean{TNLean.PEPS.EdgeMiddleRegionInjectivityComparison.edgeBlockedThreeSiteInjective}
    \leanok
    \uses{def:IsVertexInjective, def:peps_edgeBlockedThreeSiteInjective,
        def:peps_edgeMiddleRegionInjectivityComparison,
        thm:peps_edgeBlockedThreeSiteInjective_of_middle}
    Suppose the middle region $V\setminus\{u,v\}$ is injective after
    blocking, and suppose this region-injectivity assertion has been identified
    with injectivity of the corresponding edge-middle tensor family. Then a
    vertex-injective PEPS tensor gives an injective edge-blocked three-site
    chain at the edge $e=(u,v)$.
\end{theorem}
\begin{proof}\leanok
    Let $e=(u,v)$ and put $M_e=V\setminus\{u,v\}$. The comparison gives
    $\kappa(M_e)\Longrightarrow\operatorname{inj}(T_{A,M_e})$. Vertex injectivity
    supplies $\operatorname{inj}(T_{A,u})$ and $\operatorname{inj}(T_{A,v})$.
    The preceding reduction then gives
    \[
        \operatorname{inj}(T_{A,u})\wedge
        \operatorname{inj}(T_{A,M_e})\wedge
        \operatorname{inj}(T_{A,v}).
    \]
\end{proof}

\begin{theorem}[Middle tensors from all middle regions]%
    \label{thm:peps_edgeMiddleTensorInjective_all_of_region}
    \lean{TNLean.PEPS.EdgeMiddleRegionInjectivityComparison.edgeMiddleTensorInjective_all}
    \leanok
    \uses{def:peps_edgeMiddleRegionInjectivityComparison,
        def:peps_edgeMiddleTensorFamily}
    Suppose every edge-middle region $V\setminus\{u,v\}$ is injective after
    blocking, and suppose finite-region injectivity has been identified with
    injectivity of the corresponding edge-middle tensor family. Then every
    edge-middle tensor family is injective.
\end{theorem}
\begin{proof}\leanok
    For each edge $e=(u,v)$, put $M_e=V\setminus\{u,v\}$. The hypotheses give
    $\kappa(M_e)$. The comparison implication
    $\kappa(M_e)\Longrightarrow\operatorname{inj}(T_{A,M_e})$ yields
    $\operatorname{inj}(T_{A,M_e})$ for every edge.
\end{proof}

\begin{theorem}[All edge-blocked chains from middle-region injectivity]%
    \label{thm:peps_edgeBlockedThreeSiteInjective_all_of_region}
    \lean{TNLean.PEPS.EdgeMiddleRegionInjectivityComparison.edgeBlockedThreeSiteInjective_all}
    \leanok
    \uses{def:IsVertexInjective, def:peps_edgeBlockedThreeSiteInjective,
        def:peps_edgeMiddleRegionInjectivityComparison,
        thm:peps_edgeBlockedThreeSiteInjective_all_of_middle,
        thm:peps_edgeMiddleTensorInjective_all_of_region}
    Suppose every edge-middle region $V\setminus\{u,v\}$ is injective after
    blocking, and suppose finite-region injectivity has been identified with
    injectivity of the corresponding edge-middle tensor family. Then a
    vertex-injective PEPS tensor gives an injective edge-blocked three-site
    chain at every edge.
\end{theorem}
\begin{proof}\leanok
    For every edge $e=(u,v)$, put $M_e=V\setminus\{u,v\}$. The region
    hypothesis gives $\kappa(M_e)$, and the comparison gives
    $\kappa(M_e)\Longrightarrow\operatorname{inj}(T_{A,M_e})$. Vertex
    injectivity gives $\operatorname{inj}(T_{A,u})$ and
    $\operatorname{inj}(T_{A,v})$. Thus
    \[
        \operatorname{inj}(T_{A,u})\wedge
        \operatorname{inj}(T_{A,M_e})\wedge
        \operatorname{inj}(T_{A,v}).
    \]
\end{proof}

\begin{theorem}[Edge-middle tensor from singleton regions, nonempty middle case]%
    \label{thm:peps_edgeMiddleTensorInjective_of_singletons}
    \lean{TNLean.PEPS.EdgeMiddleRegionInjectivityComparison.edgeMiddleTensorInjective_of_singletons}
    \leanok
    \uses{thm:peps_finiteRegionInjectivityFromSingletons,
        def:peps_edgeMiddleRegionInjectivityComparison}
    Let $e=(u,v)$ and put $M_e=V\setminus\{u,v\}$. Assume $M_e\ne\varnothing$.
    Assume also singleton comparison, union closure for $\kappa$, and the
    comparison
    $\kappa(M_e)\Longrightarrow\operatorname{inj}(T_{A,M_e})$. If $A$ is
    vertex-injective, then the edge-middle tensor family $T_{A,M_e}$ is
    injective.
\end{theorem}
\begin{proof}\leanok
    The finite-region theorem applied to
    $M_e=\bigcup_{w\in M_e}\{w\}$ gives $\kappa(M_e)$. The edge-middle
    comparison sends this to $\operatorname{inj}(T_{A,M_e})$.
\end{proof}

\begin{theorem}[One edge-blocked chain from singleton regions, nonempty middle case]%
    \label{thm:peps_edgeBlockedThreeSiteInjective_of_singletons}
    \lean{TNLean.PEPS.EdgeMiddleRegionInjectivityComparison.edgeBlockedThreeSiteInjective_of_singletons}
    \leanok
    \uses{thm:peps_edgeMiddleTensorInjective_of_singletons,
        thm:peps_edgeBlockedThreeSiteInjective_of_middle}
    Let $e=(u,v)$ and put $M_e=V\setminus\{u,v\}$. Under the hypotheses of the
    preceding theorem, vertex injectivity of $A$ gives injectivity of the
    edge-blocked three-site chain at $e$.
\end{theorem}
\begin{proof}\leanok
    The preceding theorem gives $\operatorname{inj}(T_{A,M_e})$. The endpoint
    reduction combines this middle-tensor injectivity with vertex injectivity
    at $u$ and $v$.
\end{proof}

\begin{theorem}[All edge-blocked chains from singleton regions, nonempty middle case]%
    \label{thm:peps_edgeBlockedThreeSiteInjective_all_of_singletons}
    \lean{TNLean.PEPS.EdgeMiddleRegionInjectivityComparison.edgeBlockedThreeSiteInjective_all_of_singletons}
    \leanok
    \uses{thm:peps_edgeBlockedThreeSiteInjective_of_singletons}
    If $M_e=V\setminus\{u,v\}$ is nonempty for every edge $e=(u,v)$, then the
    preceding one-edge statement applies at every edge. Thus vertex
    injectivity of $A$ gives injectivity of every edge-blocked three-site
    chain.
\end{theorem}
\begin{proof}\leanok
    Apply the one-edge statement separately to each edge $e$.
\end{proof}

\begin{theorem}[All edge-blocked chains from singleton regions, cardinal form]%
    \label{thm:peps_edgeBlockedThreeSiteInjective_all_two_lt_card}
    \lean{TNLean.PEPS.EdgeMiddleRegionInjectivityComparison.edgeMiddleTensorInjective_all_two_lt_card}
    \lean{TNLean.PEPS.EdgeMiddleRegionInjectivityComparison.edgeBlockedThreeSiteInjective_all_two_lt_card}
    \leanok
    \uses{thm:peps_edgeMiddleVertices_nonempty_of_two_lt_card,
        thm:peps_edgeBlockedThreeSiteInjective_all_of_singletons,
        def:peps_edgeMiddleRegionInjectivityComparison}
    Assume singleton comparison, union closure for $\kappa$, the comparison
    $\kappa(M_e)\Longrightarrow\operatorname{inj}(T_{A,M_e})$, and vertex
    injectivity of $A$. If $2<|V|$, then
    $V\setminus\{u,v\}\ne\varnothing$ for every edge $e=(u,v)$, and these
    hypotheses give $\operatorname{inj}(T_{A,V\setminus\{u,v\}})$ and
    edge-blocked three-site injectivity at every edge.
\end{theorem}
\begin{proof}\leanok
    For each $e=(u,v)$, Theorem~\ref{thm:peps_edgeMiddleVertices_nonempty_of_two_lt_card}
    gives $V\setminus\{u,v\}\ne\varnothing$. Apply the preceding all-edge
    singleton theorem.
\end{proof}
